\chapter{函数与导数}
\setcounter{choicecounter}{0}

%% ---- 知识框架 ---- %%
\begin{knowledgeframe}
\begin{itemize}
	\item \textbf{函数的奇偶性}：$f(-x)=-f(x)$ 为奇函数，$f(-x)=f(x)$ 为偶函数；定义域必须关于原点对称。
	\item \textbf{函数的单调性}：$f'(x)>0$ 递增，$f'(x)<0$ 递减；含参函数需分类讨论。
	\item \textbf{对数函数的定义域}：$\lg g(x)$ 要求 $g(x)>0$；单调性分析需同时考虑定义域。
	\item \textbf{充分条件与必要条件}：函数性质判断中常用充要条件分析。
\end{itemize}
\end{knowledgeframe}

%% ---- 第一节：函数的性质 ---- %%
\section{函数的性质}

\startexercise

\subsection*{A组}

\begin{choice}{2025深圳二模}{含参函数的奇偶与单调性}
	已知函数 $f(x) = ae^x - e^{-x}$（$a$ 为常数），则
	\begin{tasks}(2)
		\task $\forall a \in \mathbb{R}$, $f(x)$ 为奇函数
		\task $\exists a \in \mathbb{R}$, $f(x)$ 为偶函数
		\task $\forall a \in \mathbb{R}$, $f(x)$ 为增函数
		\task $\exists a \in \mathbb{R}$, $f(x)$ 为减函数
	\end{tasks}
\end{choice}

\subsection*{B组}

\begin{choice}{}{对数函数的定义域与单调性}
	已知函数 $f(x) = \lg(x^2 - ax + 2)$，则"$a \ge 2$"是"函数 $f(x)$ 在 $(-\infty, 1]$ 上单调递减"的
	\begin{tasks}(2)
		\task 充分条件
		\task 充分不必要条件
		\task 必要条件
		\task 既不充分也不必要条件
	\end{tasks}
\end{choice}

\stopexercise
